%=====================================================================
%  Actividad 4D-02 : "El vinculo que no era"
%  Mecanica Clasica (posgrado) - Escuela de Fisica, UIS
%  Tema Beamer: Madrid
%
%  NOTA AL DOCENTE:
%   * El bloque \appendix del final contiene la CLAVE (dictamen R1-R8).
%     Comente esas lineas antes de distribuir el PDF a los estudiantes.
%   * Para incluir el logo institucional, descomente la linea \logo{...}
%     y coloque el archivo del logo en el mismo directorio.
%  Compilar:  pdflatex archivo.tex   (dos veces)
%=====================================================================
\documentclass[10pt,aspectratio=169]{beamer}

\usetheme{Madrid}
\usecolortheme{default}

\usepackage[utf8]{inputenc}
\usepackage[T1]{fontenc}
% Carga portable del castellano: usa spanish.ldf si esta instalado;
% si no, recurre al archivo de localizacion .ini de babel moderno.
\IfFileExists{spanish.ldf}
  {\usepackage[spanish,es-noshorthands]{babel}}
  {\usepackage[spanish,provide=*]{babel}}
\usepackage{amsmath,amssymb,bm}
\usepackage{booktabs}
\usepackage{array}
\usepackage{graphicx}

% \logo{\includegraphics[height=0.8cm]{logo-uis.png}}

\setbeamertemplate{navigation symbols}{}
\setbeamertemplate{caption}[numbered]
\setbeamercovered{transparent}

\newcommand{\rv}{\mathbf{r}}
\newcommand{\vv}{\mathbf{v}}

\title[Actividad 4D-02: vínculos]{Actividad 4D--02: «El vínculo que no era»}
\subtitle{Auditoría de un asistente IA sobre coordenadas generalizadas y vínculos}
\author[L.A. Núñez]{Luis A. Núñez}
\institute[UIS]{Escuela de Física, Facultad de Ciencias,\\ Universidad Industrial de Santander, Santander, Colombia}
\date{18 de agosto de 2026}

\begin{document}

\begin{frame}
  \titlepage
\end{frame}

%---------------------------------------------------------------------
\begin{frame}{Agenda}
  \tableofcontents
\end{frame}

%=====================================================================
\section{Objetivos y reglas de juego}
%=====================================================================

\begin{frame}{¿Qué se persigue con esta actividad?}
  \begin{itemize}[<+->]
    \item Casi todos los textos afirman que \emph{«rodar sin deslizar es un vínculo no holónomo»}.
    \item Un asistente de IA reproducirá esa afirmación \textbf{con toda confianza}, porque es estable en la literatura.
    \item Pero la clasificación de un vínculo \textbf{no depende de que aparezcan velocidades}: depende de si la forma diferencial asociada es \textbf{integrable}.
    \item Usted va a \textbf{fijar} el criterio correcto \emph{corrigiendo a la IA}, no escuchando una definición.
  \end{itemize}
  \onslide<5->{
  \begin{alertblock}{Regla del curso}
    Cada respuesta de la IA es una \textbf{hipótesis}. Nunca acepte una que no haya pasado al menos una prueba estructural y un caso límite.
  \end{alertblock}}
\end{frame}

\begin{frame}{Objetivos de aprendizaje}
  Al finalizar, usted estará en capacidad de:
  \begin{enumerate}[<+->]
    \item \textbf{Formular} los vínculos de un sistema como ecuaciones explícitas.
    \item \textbf{Clasificar} un vínculo (holónomo/anholónomo, esclerónomo/reónomo) \emph{aplicando un criterio de integrabilidad}.
    \item \textbf{Distinguir} la reducción de la dimensión del espacio de configuración de la restricción sobre las velocidades accesibles.
    \item \textbf{Construir} $\rv_i=\rv_i(q_1,\dots,q_s,t)$ y evaluar si un conjunto $\{q_j\}$ parametriza la variedad de configuración.
    \item \textbf{Auditar} una respuesta de IA con una batería explícita de pruebas.
    \item \textbf{Documentar} la interacción de modo que un tercero la reproduzca.
  \end{enumerate}
\end{frame}

\begin{frame}{Conceptos previos}
  \begin{itemize}
    \item Sistema de $N$ partículas y sus $3N$ coordenadas; cuerpo rígido libre y sus \textbf{6} grados de libertad.
    \item Vínculo holónomo $f_j(\rv_1,\dots,\rv_N,t)=0$; conteo $s=3N-k$ para $k$ vínculos holónomos \textbf{independientes}.
    \item Forma de Pfaff $\sum_j a_j\,dq_j+a_t\,dt=0$; diferencial exacta; factor integrante; derivadas cruzadas.
    \item Jacobiano, cambio de variables, singularidades de una parametrización.
    \item Cinemática de rodadura sin deslizamiento.
  \end{itemize}
\end{frame}

\begin{frame}{Política de uso de IA en esta actividad}
  \small
  \begin{table}
  \begin{tabular}{@{}p{3.1cm}p{2.9cm}p{7.4cm}@{}}
    \toprule
    \textbf{Fase} & \textbf{Uso de IA} & \textbf{Justificación} \\
    \midrule
    1 --- Delegación & \alert{Prohibido} & Identificar vínculos, contarlos y elegir las $q_j$ \emph{es} el objetivo de aprendizaje: quien no sabe hacerlo no puede verificar el resultado. \\
    2 --- Descripción & Obligatorio & Se evalúa la calidad de la especificación, no la respuesta. \\
    3 --- Discernimiento & Obligatorio & La respuesta de la IA es el objeto de estudio. \\
    4 --- Diligencia & Solo álgebra verificable y depuración de código & Toda rederivación se hace a mano y queda registrada. \\
    \bottomrule
  \end{tabular}
  \end{table}
  \vspace{0.2em}
  La Ficha de la Fase~1 se entrega con \textbf{marca de tiempo anterior} a la primera interacción registrada en la bitácora.
\end{frame}

%=====================================================================
\section{El sistema núcleo}
%=====================================================================

\begin{frame}{Núcleo común: dos casos que \emph{parecen} el mismo}
  \begin{block}{Caso A --- rodadura rectilínea}
    Disco de radio $R$ en un plano vertical \textbf{fijo}, rueda sin deslizar en línea recta.\\
    Coordenadas candidatas: $(x,\theta)$. \qquad Vínculo: $\ \dot x - R\dot\theta = 0$.
  \end{block}
  \pause
  \begin{block}{Caso B --- el disco puede cambiar de rumbo}
    Disco vertical de radio $R$ que rueda sin deslizar y \textbf{puede girar alrededor de la vertical} por el punto de contacto.\\
    Coordenadas candidatas: $(x,y,\varphi,\theta)$, con $\varphi$ el rumbo y $\theta$ el giro propio.
    \[
      \dot x = R\,\dot\theta\cos\varphi, \qquad \dot y = R\,\dot\theta\sin\varphi .
    \]
  \end{block}
  \pause
  \begin{center}
    \alert{Ambos «ruedan sin deslizar». ¿Pertenecen a la misma clase de vínculo?}
  \end{center}
\end{frame}

\begin{frame}{Sistema asignado (según el último dígito del código)}
  \small
  \begin{table}
  \begin{tabular}{@{}cp{11.5cm}@{}}
    \toprule
    \textbf{Dígito} & \textbf{Sistema} \\
    \midrule
    0, 1 & Péndulo simple plano de longitud $l$ (masa puntual, varilla rígida sin masa). \\
    2, 3 & Péndulo doble plano, $(l_1,m_1)$, $(l_2,m_2)$. \\
    4, 5 & Partícula sobre la superficie \textbf{interior} de un cono de semiángulo $\alpha$, vértice en el origen, eje $z$ vertical. \\
    6, 7 & Péndulo cuyo soporte recorre una circunferencia de radio $a$ en el plano vertical con $\omega$ constante. \\
    8, 9 & Partícula que desliza sin fricción sobre la superficie \textbf{exterior} de una esfera de radio $R$, partiendo del reposo en el polo. \\
    \bottomrule
  \end{tabular}
  \end{table}
  Todos los entregables se hacen para el \textbf{núcleo común} \emph{y} para el \textbf{sistema asignado}.
\end{frame}

%=====================================================================
\section{Fase 1 --- Delegación (sin IA)}
%=====================================================================

\begin{frame}{Fase 1 --- Delegación}
  \begin{quote}\small
    Acto cognitivo: descomponer la tarea y decidir qué se conserva.\\
    \textbf{Regla de asignación:} conserve lo que es objetivo de aprendizaje y no puede verificar; delegue lo que es verificable.
  \end{quote}
  \pause
  \textbf{1.1 Ficha de formulación} (para Caso A, Caso B y sistema asignado):
  \begin{columns}[T]
    \begin{column}{0.48\textwidth}
      \small
      \begin{itemize}
        \item Esquema con ejes, origen y símbolos definidos.
        \item Conteo bruto \textbf{con la regla usada} ($3N$ para partículas; $6$ por cuerpo rígido libre).
        \item Cada vínculo como ecuación numerada, con su forma diferencial.
        \item Argumento de independencia de los $k$ vínculos.
      \end{itemize}
    \end{column}
    \begin{column}{0.48\textwidth}
      \small
      \begin{itemize}
        \item Clasificación \textbf{con el criterio aplicado explícitamente}.
        \item $s$; y si hay vínculos no integrables: (dimensión del espacio de configuración, número de velocidades independientes).
        \item $\{q_j\}$ con su dominio.
        \item $\rv_i=\rv_i(q_1,\dots,q_s,t)$ para \emph{todas} las partículas.
        \item Número de condiciones iniciales independientes.
      \end{itemize}
    \end{column}
  \end{columns}
\end{frame}

\begin{frame}{Fase 1 --- Tabla de delegación}
  \small
  \begin{table}
  \begin{tabular}{@{}p{6.4cm}p{2.4cm}p{4.6cm}@{}}
    \toprule
    \textbf{Subtarea} & \textbf{¿Conservo / delego?} & \textbf{¿Por qué? ¿Cómo verifico?} \\
    \midrule
    Identificar los vínculos & & \\
    Clasificarlos & & \\
    Elegir las $q_j$ & & \\
    Verificar exactitud de una forma diferencial & & \\
    Buscar un factor integrante & & \\
    Álgebra de la condición de integrabilidad & & \\
    Graficar / integrar numéricamente & & \\
    Interpretar el resultado & & \\
    \bottomrule
  \end{tabular}
  \end{table}
  \begin{alertblock}{Restricción}
    Las tres primeras filas \textbf{deben} quedar del lado humano. Si su tabla las delega, la Fase~1 no cumple su objetivo.
  \end{alertblock}
\end{frame}

%=====================================================================
\section{Fase 2 --- Descripción}
%=====================================================================

\begin{frame}{Fase 2 --- Descripción: el prompt es la formulación}
  \begin{quote}\small
    Una especificación insuficiente \emph{no} produce una respuesta vaga: produce una respuesta segura al problema de libro más común.
  \end{quote}
  \pause
  \begin{enumerate}
    \item \textbf{Prompt vago}: subespecifique deliberadamente el \textbf{Caso B} (sin decir si el disco puede cambiar de rumbo, sin fijar ejes, sin pedir justificación). Registre la respuesta completa.
    \pause
    \item \textbf{Prompt preciso}: misma pregunta, fijando como mínimo\\
      sistema de referencia y origen $\cdot$ idealizaciones (rígido, delgado, contacto puntual) $\cdot$ si el plano del disco puede rotar respecto de la vertical $\cdot$ significado de cada símbolo $\cdot$ qué se pide (vínculos en forma diferencial, clasificación \textbf{y el criterio de integrabilidad usado}) $\cdot$ formato de la respuesta.
    \pause
    \item \textbf{Comparación} (media página): ¿qué supuestos rellenó la IA por su cuenta? ¿Cambian la clasificación y el conteo de $s$? ¿Alguna respuesta declara su propia incertidumbre?
  \end{enumerate}
\end{frame}

%=====================================================================
\section{Fase 3 --- Discernimiento}
%=====================================================================

\begin{frame}{Batería de verificación V1--V7}
  \small
  \begin{table}
  \begin{tabular}{@{}cp{2.7cm}p{9.4cm}@{}}
    \toprule
    & \textbf{Prueba} & \textbf{Pregunta operativa} \\
    \midrule
    V1 & Conteo & ¿La regla de partida es la correcta? ¿$s=3N-k$ se aplicó \emph{solo} a vínculos holónomos independientes? \\
    V2 & Datos iniciales & ¿El número de condiciones iniciales es $2s$? Si excede $2s$, hay vínculos no integrables y la reducción es inválida. \\
    V3 & Integrabilidad & ¿$\partial a_j/\partial q_k=\partial a_k/\partial q_j$? Si no, ¿existe factor integrante? ¿Se cumple $\omega\wedge d\omega=0$? \\
    V4 & Accesibilidad & ¿Hay un ciclo cerrado en unas coordenadas que \textbf{cambie} otra? Si sí, ninguna $F(q)=$ cte puede existir. \\
    V5 & Parametrización & ¿Las $q_j$ determinan unívocamente todos los $\rv_i$? ¿Dónde falla (jacobiano nulo, singularidades)? \\
    V6 & Tiempo y dimensiones & ¿$\partial\rv_i/\partial t\neq 0$? (esclerónomo vs.\ reónomo). ¿Homogeneidad dimensional? \\
    V7 & Dominio de validez & ¿Bilateral o unilateral? ¿Hasta qué instante o en qué región vale? \\
    \bottomrule
  \end{tabular}
  \end{table}
\end{frame}

\begin{frame}{Una regla que conviene demostrar, no memorizar}
  \begin{block}{Formas de Pfaff en dos variables}
    Una forma de Pfaff autónoma
    \[
      a(q_1,q_2)\,dq_1 + b(q_1,q_2)\,dq_2 = 0 ,
    \]
    con $a,b$ suaves y no simultáneamente nulas, \textbf{siempre} admite factor integrante local.
  \end{block}
  \pause
  \begin{center}
    \alert{Extraiga usted la consecuencia sobre el Caso A.}\\[0.4em]
    ¿Cuántas variables independientes hay en el Caso A? ¿Y en el Caso B?
  \end{center}
\end{frame}

\begin{frame}{Fase 3 --- Parte A: auditoría de la transcripción}
  Para \textbf{cada} afirmación R1--R8 del Anexo:
  \begin{enumerate}
    \item Dictamine: \textbf{correcta / incorrecta / parcialmente correcta}.
    \item Si no es correcta, cite \textbf{el fragmento exacto} que falla.
    \item Indique \textbf{qué prueba V1--V7} la detecta.
    \item Escriba la \textbf{versión corregida}, con justificación completa.
  \end{enumerate}
  \pause
  \begin{alertblock}{Advertencia}
    La transcripción contiene al menos una afirmación \textbf{correcta}. Marcarla como errónea se penaliza igual que no detectar un error: \textbf{la calibración importa tanto como la sospecha}.
  \end{alertblock}
\end{frame}

\begin{frame}{Anexo --- Transcripción a auditar (1/3)}
  \small
  \textit{Prompt utilizado:} «Explícame las coordenadas generalizadas y los vínculos del disco que rueda y de otros sistemas mecánicos típicos.»
  \vspace{0.6em}

  \textbf{R1.} «Un disco rígido que rueda es un sistema de $N=1$ cuerpo, de modo que se describe con $3N=3$ coordenadas cartesianas, sobre las que después se imponen los vínculos.»
  \vspace{0.4em}

  \textbf{R2.} «La condición de rodadura sin deslizamiento $\dot x_{cm}-R\dot\theta=0$ involucra velocidades; por lo tanto es un vínculo anholónomo y no admite integración.»
  \vspace{0.4em}

  \textbf{R3.} «Para el disco vertical que además puede orientarse sobre el plano, las condiciones $\dot x=R\dot\theta\cos\varphi$ y $\dot y=R\dot\theta\sin\varphi$ son dos vínculos, luego $s=4-2=2$ y basta elegir $q=(x,y)$ como coordenadas generalizadas.»
\end{frame}

\begin{frame}{Anexo --- Transcripción a auditar (2/3)}
  \small
  \textbf{R4.} «En la partícula sobre el cono de semiángulo $\alpha$, el vínculo $z=r\cot\alpha$ relaciona dos coordenadas entre sí y depende de la geometría del sistema, por lo que es un vínculo anholónomo, con $s=2$.»
  \vspace{0.5em}

  \textbf{R5.} «Para el péndulo doble plano con masas puntuales y varillas rígidas: $x_1^2+y_1^2-l_1^2=0$ y $(x_2-x_1)^2+(y_2-y_1)^2-l_2^2=0$ son vínculos holónomos y esclerónomos; con $N=2$ y $3N=6$, la condición de movimiento plano añade $z_1=0$ y $z_2=0$, de modo que $s=6-4=2$ y se puede tomar $q=(\theta_1,\theta_2)$.»
  \vspace{0.5em}

  \textbf{R6.} «Para el péndulo cuyo soporte recorre una circunferencia con velocidad angular constante $\omega$, el vínculo es holónomo y esclerónomo; por lo tanto la energía mecánica $E=T+V$ se conserva y $s=1$.»
\end{frame}

\begin{frame}{Anexo --- Transcripción a auditar (3/3)}
  \small
  \textbf{R7.} «Las coordenadas generalizadas deben tener dimensiones de longitud, para que las velocidades generalizadas tengan dimensiones de velocidad; por eso conviene usar la longitud de arco $s=l\theta$ en lugar del ángulo $\theta$.»
  \vspace{0.6em}

  \textbf{R8.} «Una partícula que desliza sobre la superficie exterior de una esfera lisa de radio $R$ obedece el vínculo holónomo $r-R=0$ durante todo el movimiento, de manera que $s=2$ en todo instante.»
  \vspace{1.2em}
  \pause
  \begin{center}
    \alert{Ocho afirmaciones. Al menos una es correcta. ¿Cuáles sobreviven a V1--V7?}
  \end{center}
\end{frame}

\begin{frame}{Fase 3 --- Parte B: auditoría de \emph{su propia} interacción}
  Aplique la misma batería a la respuesta obtenida con su prompt preciso. Identifique:
  \begin{itemize}
    \item errores;
    \item afirmaciones no verificables;
    \item afirmaciones correctas pero mal justificadas;
    \item \textbf{silencios}: lo que la IA omitió y su Ficha de la Fase~1 sí contemplaba.
  \end{itemize}
  \pause
  \begin{block}{Interacción de retorno (obligatoria)}
    Presente a la IA \textbf{una} de sus correcciones y registre su reacción.
    ¿Cede sin argumentos? ¿Defiende su posición? ¿Introduce un error nuevo al corregir?\\
    Comente en tres líneas qué le dice esto sobre \textbf{el valor probatorio de la conformidad de un modelo}.
  \end{block}
\end{frame}

%=====================================================================
\section{Fase 4 --- Diligencia}
%=====================================================================

\begin{frame}{Fase 4 --- Diligencia}
  \begin{enumerate}
    \item \textbf{Rederivación independiente y a mano} de:
      \begin{itemize}\small
        \item la integración explícita del vínculo del \textbf{Caso A}, con la constante interpretada físicamente;
        \item la demostración de que el sistema de vínculos del \textbf{Caso B} no es integrable;
        \item el conteo de grados de libertad y de datos iniciales en ambos casos;
        \item el punto o condición donde el vínculo de su sistema asignado deja de ser válido, si aplica.
      \end{itemize}
    \pause
    \item \textbf{Bitácora 4D}: \emph{prompt $\cdot$ respuesta relevante $\cdot$ aceptado/rechazado/corregido $\cdot$ cómo se verificó (prueba V) $\cdot$ error o limitación detectada $\cdot$ razonamiento final propio}.
    \pause
    \item \textbf{Declaración de uso de IA}: herramienta y versión, fecha, fases de uso, y la frase\\
      \emph{«Verifiqué de forma independiente los resultados sobre los que construí; los errores que subsisten son míos».}
    \pause
    \item \textbf{Reproducibilidad}: código documentado, tolerancias y semillas declaradas.
  \end{enumerate}
\end{frame}

%=====================================================================
\section{Extensión computacional}
%=====================================================================

\begin{frame}{Extensión computacional (opcional, recomendada)}
  Demuestre \textbf{numéricamente} la no integrabilidad del Caso B por accesibilidad (V4).
  \begin{block}{Tarea}
    Integre $\ \dot x = R\dot\theta\cos\varphi$, $\ \dot y = R\dot\theta\sin\varphi$, imponiendo que el punto de contacto recorra una circunferencia de radio $R_c$ con rapidez $v$ constante
    ($\dot\theta=v/R$, $\dot\varphi=v/R_c$) durante $T=2\pi R_c/v$.
  \end{block}
  \pause
  \textbf{Reporte:}
  \begin{enumerate}\small
    \item Valores finales de $(x,y,\varphi,\theta)$ frente a los iniciales.
    \item El argumento: si existiera $F(x,y,\varphi,\theta)=$ cte, un ciclo que devuelve $x,y,\varphi$ debería devolver $\theta$. Verifique $\Delta\theta = 2\pi R_c/R$ y concluya.
    \item Control de convergencia: dos tolerancias distintas, error en $\Delta\theta$ frente al valor analítico.
    \item Repita con $R_c/R = 1,\,3,\,10$ y compruebe la ley.
  \end{enumerate}
  \pause
  \begin{center}\small
    Las ecuaciones se derivan \textbf{antes} de implementarlas. El código puede depurarse con IA;
    \alert{la interpretación de la salida no se delega.}
  \end{center}
\end{frame}

%=====================================================================
\section{Entregables, defensa y rúbrica}
%=====================================================================

\begin{frame}{Entregables}
  \small
  \begin{enumerate}
    \item Ficha de formulación (Fase 1), con marca de tiempo previa a la interacción con IA.
    \item Tabla de delegación.
    \item Anexo de prompts: vago, preciso y respuestas completas \emph{sin editar}.
    \item Auditoría R1--R8 en tabla: dictamen $\cdot$ fragmento que falla $\cdot$ prueba V $\cdot$ corrección justificada.
    \item Auditoría de la interacción propia + comentario de la interacción de retorno.
    \item Rederivaciones a mano (manuscritas o en \LaTeX).
    \item Bitácora 4D (seis columnas).
    \item Declaración de uso de IA.
    \item (Opcional) Código y figuras de la extensión computacional.
  \end{enumerate}
  \vspace{0.3em}
  \textbf{Formato:} máximo 8 páginas sin contar anexos de transcripciones. Notación consistente con la del curso.
\end{frame}

\begin{frame}{Defensa oral (5 minutos)}
  El docente elegirá dos preguntas de variación:
  \begin{itemize}\small
    \item Cambio el disco del Caso A por una esfera que rueda sin deslizar sobre el plano. ¿Cambia su conclusión? ¿Por qué?
    \item Muéstreme en la forma diferencial dónde exactamente falla la exactitud en el Caso B.
    \item ¿Cuántas condiciones iniciales necesito en el Caso B? ¿Y si el disco pudiera además deslizar?
    \item Si el disco del Caso A puede despegarse del piso, ¿qué tipo de vínculo aparece?
    \item Su compañero dice que el Caso B tiene 2 grados de libertad porque hay 2 vínculos. Convénzalo \emph{sin escribir ecuaciones}.
    \item ¿Puede el vínculo de su sistema asignado ser reónomo si elijo otro sistema de referencia? Argumente.
  \end{itemize}
\end{frame}

\begin{frame}{Rúbrica (100 puntos)}
  \scriptsize
  \begin{table}
  \begin{tabular}{@{}p{4.0cm}cp{8.2cm}@{}}
    \toprule
    \textbf{Criterio} & \textbf{Peso} & \textbf{Nivel Ejemplar (4/4)} \\
    \midrule
    1. Formulación física propia & 20 & Vínculos, conteo, transformaciones y datos iniciales correctos en ambos casos y en el sistema asignado; dominios y singularidades explícitos. \\
    2. Delegación & 10 & Conserva identificación, clasificación y elección de $q$; enuncia \emph{cómo} verificaría cada delegación. \\
    3. Descripción & 15 & El prompt preciso fija marco, idealizaciones, símbolos, criterio y formato; la comparación identifica los supuestos rellenados por la IA. \\
    4. Discernimiento --- Anexo & 25 & Dictamina correctamente los 8 ítems (incluido R5); cita el fragmento, asigna la prueba V y corrige con justificación completa. \\
    5. Discernimiento --- interacción propia & 10 & Aplica $\geq 4$ pruebas V con resultado explícito; comenta críticamente la interacción de retorno. \\
    6. Diligencia & 15 & Rederivaciones completas y correctas; bitácora de seis columnas; declaración firmada; código reproducible. \\
    7. Comunicación y notación & 5 & Notación consistente; símbolos definidos; extensión respetada. \\
    \bottomrule
  \end{tabular}
  \end{table}
  \vspace{-0.2em}
  \textbf{Defensa oral:} $\pm 10$. \quad
  \textbf{Integridad:} una auditoría correcta cuya bitácora no permita reconstruir el camino no supera el nivel 2 en el criterio 6.
\end{frame}

%=====================================================================
\section{Para la discusión}
%=====================================================================

\begin{frame}{Para la discusión}
  \begin{enumerate}[<+->]
    \item ¿Por qué la heurística «hay velocidades $\Rightarrow$ es anholónomo» sobrevive en tantos textos? ¿Qué la hace tan estable en la salida de un asistente de IA?
    \item Si un vínculo no integrable \emph{no} reduce la dimensión del espacio de configuración, ¿qué es exactamente lo que reduce?
    \item ¿Qué prueba de V1--V7 habría bastado, por sí sola, para detectar el error de R3?
    \item ¿Puede una respuesta ser \emph{correcta} y aun así inaceptable como entregable de este curso?
    \item Si la IA acepta su corrección de inmediato, ¿ha aprendido usted algo sobre si tenía razón?
  \end{enumerate}
\end{frame}

\begin{frame}{Recapitulando}
  \begin{enumerate}
    \item \textbf{Criterio, no mnemotecnia:} la presencia de velocidades es condición \emph{necesaria}, no suficiente, para la no holonomía.
    \item \textbf{Dos reducciones distintas:} los vínculos holónomos reducen la dimensión del espacio de configuración; los no integrables restringen las velocidades accesibles.
    \item \textbf{Las 4D en cadena:} Delegación $\to$ Descripción $\to$ Discernimiento $\to$ Diligencia, unidas por la compuerta de verificación.
    \item \textbf{La calibración es parte de la competencia:} detectar todo lo falso \emph{y} no marcar lo verdadero.
  \end{enumerate}
\end{frame}

%=====================================================================
%  CLAVE DEL DOCENTE --- comentar antes de distribuir a estudiantes
%=====================================================================
\appendix

\begin{frame}{\textsc{Clave del docente} --- dictamen R1--R8}
  \small
  \begin{table}
  \begin{tabular}{@{}ccp{8.6cm}@{}}
    \toprule
    \textbf{Ítem} & \textbf{Dictamen} & \textbf{Prueba / motivo} \\
    \midrule
    R1 & Incorrecta & V1: $3N$ vale para partículas; un cuerpo rígido libre tiene 6 g.d.l.; el Caso B usa $(x,y,\varphi,\theta)$. \\
    R2 & \alert{Incorrecta} & V3: $dx-R\,d\theta$ es exacta $\Rightarrow x-R\theta=$ cte. Holónomo y esclerónomo, $s=1$. \\
    R3 & Incorrecta (dos errores) & V2--V5: no integrable \emph{y} los vínculos no integrables no reducen el espacio de configuración. \\
    R4 & Parcialmente correcta & V3: $z=r\cot\alpha$ es holónomo y esclerónomo; $s=2$ sí es correcto. \\
    R5 & \textbf{Correcta} & Control de falsos positivos: $s=6-4=2$, $q=(\theta_1,\theta_2)$. \\
    R6 & Parcialmente correcta & V6: holónomo pero \textbf{reónomo}; $E=T+V$ no se conserva; $s=1$ correcto. \\
    R7 & Incorrecta & V6: las $q_j$ pueden tener dimensiones distintas o ninguna; $[p_jq_j]=$ acción. \\
    R8 & Incorrecta & V7: vínculo \textbf{unilateral} $r\ge R$; se pierde en $\cos\theta_c=2/3$. \\
    \bottomrule
  \end{tabular}
  \end{table}
\end{frame}

\begin{frame}{\textsc{Clave del docente} --- los dos puntos centrales}
  \small
  \begin{block}{Caso A es holónomo (R2)}
    $\omega = dx - R\,d\theta$, con $\partial(1)/\partial\theta = \partial(-R)/\partial x = 0$: exacta.
    $\Rightarrow x-R\theta=$ cte (posición del contacto cuando $\theta=0$); $s = 2-1 = 1$, $q=\theta$.
    Refuerzo: en dos variables \emph{siempre} hay factor integrante $\Rightarrow$ la no holonomía exige al menos tres variables.
  \end{block}
  \begin{block}{Caso B no es integrable (R3)}
    Con $\omega_1 = dx - R\cos\varphi\,d\theta$: \ $d\omega_1 = R\sin\varphi\,d\varphi\wedge d\theta$, luego
    \[\omega_1\wedge d\omega_1 = R\sin\varphi\;dx\wedge d\varphi\wedge d\theta \neq 0 \quad\text{(Frobenius falla).}\]
    Configuración: 4 coordenadas accesibles; velocidades independientes: 2; datos iniciales: $4+2=6 \neq 2s=4$.
    Accesibilidad: rodando por una circunferencia de radio $R_c$ el disco regresa a $(x,y)$ y a $\varphi$, pero $\Delta\theta = 2\pi R_c/R$.
    \emph{(Verificado numéricamente: $R=1$, $R_c=3$ $\Rightarrow$ $x,y\sim10^{-13}$, $\varphi/2\pi = 1.000$, $\theta = 18.8496 = 6\pi$.)}
  \end{block}
\end{frame}

\begin{frame}{\textsc{Clave del docente} --- errores frecuentes e intervención}
  \scriptsize
  \begin{table}
  \begin{tabular}{@{}p{4.6cm}p{4.2cm}p{5.0cm}@{}}
    \toprule
    \textbf{Error} & \textbf{Señal en el entregable} & \textbf{Intervención} \\
    \midrule
    Clasificar por presencia de velocidades & R2 marcado anholónomo «porque hay $\dot x$» & Exigir el cálculo de derivadas cruzadas antes del dictamen. \\
    Aplicar $s=3N-k$ a vínculos no integrables & $s=2$ en el Caso B & Pedir el conteo de condiciones iniciales (V2). \\
    Confundir cuerpo rígido con partícula & $3N=3$ en R1 & Preguntar cuántos ángulos orientan el disco. \\
    Elegir $q$ que no parametrizan & $q=(x,y)$ en el Caso B & Pedir escribir $\rv_i(q,t)$: el intento falla solo. \\
    Confundir esclerónomo con «el movimiento no depende del tiempo» & R6 clasificado esclerónomo & Pedir $\partial\rv/\partial t$ explícita. \\
    Marcar R5 como error & Falso positivo & Discutir el costo epistémico de la sospecha indiscriminada. \\
    Aceptar que la IA «acepte» la corrección & Bitácora sin verificación posterior & La conformidad del modelo no es evidencia. \\
    \bottomrule
  \end{tabular}
  \end{table}
\end{frame}

\end{document}
